\paragraph{Closing the higher-seven ramified common-commutant branch.}
The separately reviewed corrected v2 argument proves scalar joint commutant
of the actual pair for every odd $N$ divisible by $245$.
Write $N=5^e n$, $5\nmid n$, $49\mid n$. A normalized integral lattice of
actual commuters, rather than the naive reduced pair, has leading coefficients
controlled by squared Gaussian binomials. Uniform graded valuation bounds
remove all lower/diagonal corrections. The elementary $q$-Lucas identity
then identifies a mixed-digit upper-Toeplitz algebra. An invertible
highest-to-lowest high-digit Fourier block reduces its leading coefficient
to a shifted prime-to-five base commutant; induction removes all remaining
high-digit coefficients. Saturation of the actual integral commuter lattice
lifts scalarity to characteristic zero.

The shifted-base step uses the already reviewed polynomial identity with
independently variable pole parameter. Its new finite checks retain six
literal effective-gauge witnesses and six distinct unit-upper-gauge witnesses
at dimension $49$. The original finite-witness naming discrepancy is explicitly
corrected by geometric-gauge transfer; the old rank-$42$ endpoint attempt is
not counted as a successful certificate. Finite bases and dictionaries alone
do not replace the uniform lattice proof.

Combining the retained seven-primary components with the separately attributed
FM-008 result yields scalar complex joint commutant of $(A_N,M_N)$ for every
odd $N\ge3$ except the genuine exception $N=7$. Each component retains its
own review and source boundary. The literal-$U_N$ corollary is restricted to
algebraic-entry commuters via the existing coefficient-separation argument.
No scalar full complex commutant of a single $U_N$, seven-divisible spectral
simplicity, quantitative spectral gap, signed trace comparison or BGS follows.


At the exceptional dimension seven, a separate exact characteristic-zero
check gives $T^2=-7$ and $G^3=-7TI=T^3I$, where $T=\operatorname{tr}G$
and $G_{xy}=\zeta_7^{-(3x^2-2xy+3y^2)}$.
Thus $H=G/T$ has $H^3=I$ and trace one; its cube-root multiplicities are
$(3,2,2)$. Since $A_7=(9/10)I$, the complex pair commutant has dimension
$3^2+2^2+2^2=17$. This is a characteristic-zero identity, with no residue
nullity lifting.
